\begin{lemma}[CCO inheritance by invariant equivalence classes]
\label{lem:rh-class-inherits}
Let $R\subseteq\mathcal D$ be a global equivalence support and $C$
any class of $E_R$. In its induced language with diversity palette
$R$, the structure on $C$ is transitive, homogeneous, triangle-defined,
and satisfies semantic homomorphic CCO (whether $C$ is finite or
infinite).
\end{lemma}
\begin{proof}
The diversity colours occurring in $C$ are exactly $R$. The inclusion
into $R$ holds because $C$ is an $E_R$-class. Conversely, for any
$r\in R$, take an actual $r$-edge of $M$ and use transitivity to send
its first endpoint to a fixed point of $C$; its second endpoint must
then also belong to $C$. The same transitivity argument sends any
point of $C$ to any other while preserving the whole class.

Any nonempty finite partial isomorphism within $C$ extends by
homogeneity of $M$ to an automorphism $g$ of $M$. Since it maps at
least one point of $C$ to another point of $C$, invariance of $E_R$
forces $g(C)=C$. The empty partial map extends by the identity.
This proves homogeneity of $C$.

A finite complete $R$-coloured graph all of whose triangles are legal
in $C$ has all its triangles legal in $M$. Triangle-definedness of
$M$ realizes it there. Every two realised vertices are $E_R$-related,
so all lie in a single $E_R$-class. One ambient automorphism moves
that class to $C$. Thus the induced age of $C$ is triangle-defined.

Finally, let $P$ be a finite partial $R$-coloured diagram and suppose
each specified simple cycle admits a homomorphism into $C$.
The same cycles map to $M$, so raw semantic homomorphic CCO yields
$f:P\to M$. Each connected component of the graph of specified
edges is mapped entirely into one $E_R$-class: all its edge colours
belong to $R$. By transitivity an ambient automorphism can move
that image class to $C$. Apply a potentially different automorphism
to each component and combine the maps. They agree with every
specified edge because there are no specified edges between distinct
components. This gives a homomorphism $P\to C$, proving semantic CCO.
\end{proof}

\begin{lemma}[Invariant equivalences below a class are exactly its internal equivalences]
\label{lem:rh-invariant-class-interval}
Let $E=E_R$ be ANY original
$\Aut(M)$-invariant equivalence
and let $C$ be ONE actual
$E$-class. Restriction to $C$
gives a lattice isomorphism
\[
 \{H:H\text{ an original
 invariant equivalence on }M,\
                 \Id\le H\le E\}
 \ \xrightarrow[\ \cong\ ]{\ H\mapsto H|_C\ }\
 \{\text{invariant equivalences
                    of the induced }C\}.
 \tag{CEI}\label{eq:rh-class-interval}
\]
Here the right side means
invariance under the FULL
$\Aut(C)$, not merely under a
preselected subgroup of
$\Aut(M)$.
This statement needs original
ONE-SORT homogeneous COMPLETE
BINARY atomicity and
transitivity, but neither CCO
nor SAP nor finiteness of $C$.

In particular, any internal
invariant equivalence on a
FINITE real one-point fibre $Z$
has a UNIQUE original global
invariant extension EXACT for
$Z$, once the finite-fibre
globalisation theorem has shown
that $Z$ is ONE class of its
global $E_T$. Thus the choice
of a saturated internal Hasse
chain in that theorem does not
affect the final extension.
\end{lemma}
\begin{proof}
By \cref{lem:rh-class-inherits},
the induced structure on $C$
is homogeneous and transitive
with precisely the actual old
diversity atom names $R$.
For every $r\in R$, an old
$r$-edge occurs in $C$:
take any old $r$-edge and move
its first endpoint to a chosen
point of $C$ by original
transitivity. Its other
endpoint is then in $C$
because $r\subseteq E$.
Moreover each old $r$-atom
restricted to $C$ is ONE
$\Aut(C)$-orbital: two
ordered $r$-pairs in $C$
give a two-point partial
isomorphism of $C$, which
extends by its homogeneity.

Let $V$ be ANY internal
$\Aut(C)$-invariant
equivalence. Its diversity
support is therefore a union
$T\subseteq R$ of ENTIRE
original named fine atom
sets, with actual logical
equality separate. Extend
it to the binary relation
$H_T=\Id_M\cup\bigcup T$
on all of $M$. This is
clearly globally invariant,
symmetric and reflexive,
and is contained in $E$.
For transitivity, if
$xH_Ty$ and $yH_Tz$,
all three points lie in ONE
original $E$-class.
Choose one original
automorphism carrying $x$
into $C$. Since it
preserves $E$, it carries
$y,z$ into that SAME $C$.
Both corresponding edges
belong to $V$ there,
so internal transitivity
forces the edge from
the image of $x$ to the
image of $z$ to belong
to $V$. Its literal old
colour is in $T$ or is
logical equality. Hence
$xH_Tz$, and $H_T$ is
an original invariant
equivalence below $E$.
By construction
$H_T|_C=V$.

Conversely every original
invariant $H\le E$
restricts to an internal
invariant equivalence,
and EVERY old atom of
$H$ occurs in $C$, so its
global diversity support
is uniquely recovered
from $H|_C$. The two
maps are therefore
inverse order
isomorphisms of the two
finite equivalence
lattices. An order
isomorphism automatically
preserves meets and joins;
no assertion of
DISTRIBUTIVITY is needed.

Finally, for a finite
one-point fibre $Z$,
\cref{thm:rh-globalise}
establishes $Z=[x]_{E_T}$.
Every $H\le E_T$ has
all its classes contained
in $Z$ whenever they meet
$Z$, so the unique lift
from the preceding
isomorphism is EXACT
for $Z$. This conclusion
is downstream of
globalisation and not
a substitute for proving
the global existence of
$E_T$.
\end{proof}

\begin{corollary}[Relative index two throughout finite original invariant intervals]
\label{cor:rh-finite-invariant-interval-index-two}
Let $E$ be an original
invariant equivalence
whose classes are FINITE.
Then EVERY Hasse cover
$H<K$ in the interval
$[\Id,E]$ of ORIGINAL
invariant equivalences
has relative index
\[
                          [K:H]=2 .
\]
Consequently every
$E$-class has size a
power of two, and the
interval $[\Id,E]$
need NOT be
distributive.
\end{corollary}
\begin{proof}
Fix an $E$-class $C$.
By \cref{lem:rh-class-inherits}
it is a finite transitive
homogeneous symmetric
structure satisfying
semantic homomorphic CCO.
The lattice isomorphism
\cref{lem:rh-invariant-class-interval}
carries any original cover
$H<K\le E$ to a Hasse cover
of invariant equivalences
on $C$. By
\cref{cor:hasse} that
internal cover has
relative index two.
The original relative
index is the SAME number
because each original
$K$-class involved is
contained in $C$ and
all its $H$-classes are
as well. Saturating a
chain from $\Id$ to $E$
gives the two-power
class-size assertion.
Nothing in the
argument forces
distributivity: finite
elementary abelian
regular $2$-groups
already have
nondistributive
subgroup intervals.
\end{proof}

\begin{example}[A nondistributive finite invariant interval in a genuine CCO source]
\label{ex:rh-finite-invariant-M3}
Take the one-sort original
homogeneous structure
\[
          M=I\times\mathbb F_2^2,
\]
where $I$ is countably
infinite. For distinct
$(i,u),(j,v)$ name one
symmetric old atom $b$
when $i\ne j$, and
when $i=j$ name the
three distinct atoms
$a_w$ for
$w=u+v\in\mathbb F_2^2
                  \setminus\{0\}$.
Its automorphism group
permutes the $I$-blocks
freely and independently
translates each
$\mathbb F_2^2$-block.
The finite induced age
has AP and semantic
homomorphic CCO:
the subgraph of specified
$a_w$ edges is a
group-labelled gain graph.
Its colour requirements
are soluble within each
connected component iff
the sum of labels around
every specified simple
cycle is zero.
A specified $b$-edge
between vertices in the
SAME such component
would form a forbidden
simple cycle together with
an $a_w$-path between
its endpoints.
Otherwise send different
$a_w$-components to
different $I$-blocks.
This gives one simultaneous
homomorphic completion.

Let $E$ be the original
invariant SAME-$I$-block
equivalence. Its classes
have four points and
the interval $[\Id,E]$
has exactly the THREE
different subequivalences
\[
        \Id\cup a_w
           \qquad
       (0\ne w\in\mathbb F_2^2)
\]
between $\Id$ and $E$.
They are pairwise
incomparable and meet
in $\Id$, while any two
join to $E$. Hence
$[\Id,E]$ is the
nondistributive
five-element diamond
$M_3$, although EVERY
Hasse cover has relative
index two. This is an
interval in the FULL
ORIGINAL invariant
equivalence lattice,
not an assertion that
all these relations
belong to the much
smaller FULL RAW exact
equivalence lattice.
\end{example}

\begin{corollary}[Two-power size of finite invariant classes]
\label{cor:rh-class-2power}
If $C$ is a finite class of an invariant equivalence $E_R$, then
\[
             |C|=2^m\quad\text{for some }0\le m\le |R|.
\]
\end{corollary}
\begin{proof}
By Lemma~\ref{lem:rh-class-inherits}, $C$ is finite, transitive,
and has semantic homomorphic CCO. In a saturated chain from equality
to the universal equivalence on $C$, every Hasse cover has index two
by Corollary~\ref{cor:hasse}. Thus $|C|=2^m$, where $m$ is the number
of covers. Since $C$ is homogeneous, its nonidentity pair orbitals
are precisely the colours in $R$. Each strict inclusion of invariant
equivalences adds at least one such orbital. Hence $m\le |R|$.
\end{proof}

